\chapter{Conclusiones}
\label{cap:conclusiones}

El presente trabajo tuvo como objetivo el diseño de una arquitectura de referencia para
Autoridades de Certificación reguladas bajo el marco legal ecuatoriano, orientada a superar
las limitaciones de rigidez, acoplamiento fuerte y baja escalabilidad de los sistemas PKI
tradicionales. El proceso comprendió cuatro fases: análisis de requisitos normativos y
funcionales, diseño de la arquitectura con el modelo C4 bajo TOGAF ADM, implementación de
una prueba de concepto del módulo de auditoría con Hyperledger Fabric, y evaluación mediante
el experimento GQM y el análisis ATAM. Los hallazgos de cada fase se sintetizan a continuación.

El análisis de los estándares PKI (X.509, RFC~5280, RFC~6960, RFC~3161) y del marco
regulatorio ecuatoriano encabezado por la Resolución ARCOTEL-2024-0176 permitió derivar
una cadena de requisitos que parte de los conductores de negocio y desciende hasta
escenarios de atributo de calidad verificables. Los seis business drivers más destacados
(Tabla~\ref{tab:bd_nfr_core}) se descompusieron en los doce NFRs y once QAS más
relevantes para el diseño (Tablas~\ref{tab:nfr_core} y~\ref{tab:qas_core}). El ejercicio confirmó que
la regulación vigente impone restricciones determinantes sobre la arquitectura: la exigencia
de respuesta OCSP inmediata tras revocación y el almacenamiento de claves exclusivamente en
HSM certificado hacen inviables los estilos de arquitectura con consistencia eventual o sin
aislamiento criptográfico de hardware. Estas restricciones se formalizaron en la
Tabla~\ref{tab:restricciones_arq} y constituyeron el filtro principal para la selección del
estilo arquitectónico. 

Con esos requisitos como marco, la evaluación sistemática de nueve estilos arquitectónicos
frente a diez criterios ponderados (Figura~\ref{fig:estilos_chart}) derivó en la selección
del estilo compuesto EA-10: Service-Based como base, complementado con asincronía selectiva
tomada de Event-Driven Architecture y extensibilidad basada en el patrón Microkernel.
Service-Based fue el único estilo con cero evaluaciones negativas, satisfaciendo los seis
criterios de peso alto y presentando solo resultado parcial en gestión de carga; esta
limitación se subsanó con la incorporación del broker de mensajería asíncrono (ADR-27,
ADR-31). La arquitectura resultante se materializa en diez contenedores de grano grueso
(Figura~\ref{fig:c4_n2} y Tabla~\ref{tab:contenedores}), cada uno con responsabilidad
delimitada por al menos una decisión arquitectónica documentada en el ADR Log. Los treinta
y nueve ADRs agrupados en seis categorías (Tabla~\ref{tab:adrs_core} y Anexo~G) formalizan
la justificación de cada decisión de diseño, vinculándola con los NFRs y QAS que la motivan
y con las métricas que la validan. Entre estas, la restricción de acoplamiento eferente
Ce~$\leq$~5 por contenedor (ADR-33), verificable directamente en los diagramas C4 de
Nivel~3, constituye una métrica objetiva que permite auditar el grado de independencia entre
contenedores sin herramientas de análisis estático externas.

Sobre esa base, la implementación del prototipo materializó esas decisiones en el módulo M-05 Immutable
Log Service. El módulo se descompone en 19 clases organizadas en tres grupos
funcionales: los controladores de entrada para escritura y consulta, el núcleo de
orquestación y persistencia centrado en \texttt{LogEntryWriter}, y la infraestructura de
mensajería asíncrona que implementa el contrato \texttt{IMessageBroker} del canal M-10. Las
decisiones de diseño más determinantes de la implementación son tres. Primero,
\texttt{HashChainEngine} opera siempre activo e independiente de la configuración de Fabric,
garantizando la integridad local del log por construcción y no por convención. Segundo,
\texttt{FabricRouteConfig} separa la política de notarización del código: añadir un nuevo
origen de eventos al ledger distribuido requiere solo una entrada en
\texttt{fabric-routes.yaml}, sin modificar ninguna clase. Tercero, \texttt{HsmSigningBridge}
abstrae la interfaz con el hardware criptográfico de M-04, confinando cualquier cambio de
algoritmo o proveedor a ese contenedor sin afectar a M-05.

La validez de este prototipo fue evaluada con el
experimento GQM descrito en el \hyperref[anx:poc_entorno]{Anexo~H}. Los resultados
confirman tres propiedades estructurales del diseño. En primer lugar, la selectividad de la
notarización es correcta: solo los eventos de alta relevancia regulatoria generan
transacciones en Hyperledger Fabric, mientras que el resto se registra exclusivamente en
el log local, sin que esta distinción comprometa la integridad del hash chain ni la
completitud de la trazabilidad por número de serie (Tabla~\ref{tab:poc_funcionales}).
En segundo lugar, la integración asíncrona con Fabric no degrada el rendimiento del canal de
mensajería bajo ninguno de los cinco escenarios de carga evaluados, lo que confirma la
viabilidad del patrón \textit{fire-and-forget} adoptado en ADR-29
(Tabla~\ref{tab:poc_estres}); la latencia de operaciones criptográficas en hardware HSM
físico queda fuera del alcance de esta PoC por restricciones económicas de adquisición.
En tercer lugar, la protección del log opera en capas
independientes: el hash chain detecta cualquier alteración retroactiva sin necesidad de
consultar el ledger, y el control de acceso rechaza solicitudes no autorizadas antes de
que alcancen el mecanismo criptográfico (Tabla~\ref{tab:poc_seguridad}). La combinación
de estas tres propiedades permite reconstruir el ciclo de vida completo de un certificado
y auditar operaciones en lote (Tabla~\ref{tab:poc_forense}).

De igual manera, el análisis ATAM examinó la arquitectura de referencia en su conjunto
(Tablas~\ref{tab:atam_utility_tree},~\ref{tab:atam_st_risk} y~\ref{tab:atam_st_nonrisk})
e identificó dos categorías de hallazgos diferenciadas. En los puntos donde la arquitectura
fija la decisión de forma estructural (estabilidad de las interfaces de extensibilidad,
aislamiento de flujos críticos respecto a la mensajería asíncrona, ausencia de estado
compartido entre instancias de servicio), las consecuencias son completamente comprendidas y
aceptables: la arquitectura de referencia garantiza esas propiedades por construcción, con
independencia del adoptante. En contraste, las brechas más relevantes
(Tabla~\ref{tab:atam_risks}) se concentran en decisiones que la arquitectura necesariamente
delega al adoptante: la completitud del registro de auditoría ante rutas de invocación que
puedan eludir el mecanismo transversal de intercepción, y la efectividad del aislamiento de
datos personales a nivel de control de acceso en la base de datos. En ambos casos, la arquitectura de referencia prescribe el
mecanismo correcto pero no puede verificar por construcción que la implementación concreta
del adoptante lo respete; esta distinción entre garantía arquitectónica y garantía de
implementación delimita el alcance de las propiedades que una arquitectura de referencia
puede ofrecer. La incertidumbre sobre la latencia de Fabric como implementación del ledger
de auditoría, identificada por el ATAM, queda resuelta por los resultados del experimento (Tabla~\ref{tab:poc_estres}): la
latencia de confirmación observada se mantuvo dentro del umbral establecido en todos los
escenarios evaluados.

Como resultado final, la arquitectura de referencia resultante constituye una especificación formal, trazable y
portable para Autoridades de Certificación que operen bajo el marco regulatorio ecuatoriano.
A diferencia de adaptaciones de plataformas internacionales sin ajuste normativo, esta
propuesta parte de los mandatos legales como restricciones de primer nivel
(Tabla~\ref{tab:restricciones_arq}) y los transforma en decisiones de diseño auditables
mediante el catálogo de ADRs (Anexo~G), los principios arquitectónicos
(Tabla~\ref{tab:principios_inline}) y los diagramas C4 en cuatro niveles de abstracción
(Figuras~\ref{fig:c4_n1} a~\ref{fig:c4_n4}). Una CA adoptante puede verificar cada decisión
contra su contexto regulatorio, sustituir los componentes de terceros por alternativas
conformes a los estándares prescritos (PKCS\#11, RFC~6960, RFC~3161) e incorporar nuevos
perfiles de certificado normativos sin rediseñar el motor de emisión. Esta capacidad de
adopción graduada, posibilitada por la separación entre estándares prescritos e
implementaciones elegibles, constituye la contribución estructural de la propuesta al
ecosistema ecuatoriano regulado de infraestructura de clave pública.

Finalmente, como líneas de trabajo futuro, la validación experimental puede extenderse a
los nueve módulos restantes de la arquitectura siguiendo el mismo proceso GQM aplicado al
módulo de auditoría, con especial atención al motor de emisión y al servicio de validación
OCSP, cuyos atributos de latencia y disponibilidad bajo carga determinan el cumplimiento
de los requisitos de mayor criticidad regulatoria. La integración con HSM físico
FIPS~140-2~L3 permitiría además completar la validación de rendimiento de extremo a
extremo, condicionada al coste de adquisición del hardware. El modelado e implementación 
completa de CA Core Service deberá incorporar además un
control estricto del ciclo de vida de variables sensibles en memoria volátil, línea de
trabajo identificada en R-NR-J. En el plano normativo,
si la regulación ecuatoriana evoluciona para permitir la generación del par de claves en
dispositivo del suscriptor, CA Core Service debería adaptarse para aceptar
solicitudes de firma de certificado (CSR) en lugar de desencadenar la generación
centralizada, lo que permitiría fortalecer la garantía de no repudio en sentido estricto.
Por último, la aplicación de la arquitectura en un proceso formal de acreditación ante
ARCOTEL proporcionaría evidencia empírica sobre la completitud de la cobertura regulatoria
en condiciones de producción.